\documentclass[main]{subfiles}
\usepackage[sols]{workbook}
\numbering{ws1-1}

\begin{document}

\ifSubfilesClassLoaded{%
  \chaptermark{\chapterone}
}{}

\section{Density and Approximation} \label{ws1-1}

\begin{enumerate}
\item  \label{constant-density-warmup}

  \note{This problem is just to get them thinking in the right direction.}  

\begin{problem}[0.5in]
    According to the US Census, the population density in Cambridge in the year 2020 was 18,500 people 
    per square mile. The area of Cambridge is about 6.4 square miles. What was the population of
Cambridge in the year 2020?
  \end{problem}

  \begin{solution}
    The population can be calculated easily:
    \begin{align*}
      \textrm{population} &= \textrm{(population density)} \cdot \textrm{(area
        of city)} \\ 
      & = (18,500 \textrm{ people per square mile}) \cdot
      (6.4 \textrm{ square miles}) \\
      & = \boxed{118,400 \textrm{ people}}
    \end{align*}
  \end{solution}

  
\item  \label{density-rectangle-1}

  \note{The exact population is 352,000 people.  Using $n = 4$ with ``right" endpoints gives the
    approximation 300,000 (using a ``left" endpoint approximation gives
    396,000.)}
  
  \begin{problem}[\vfill]
    A seaside village, Playa del Carmen, is in the shape of a rectangle
    with sides of 4 miles and 6 miles.  The sea lies along a 6-mile long
    side.  People prefer to live near the water, so the density of people
    is given by $\rho(x) = 1000 (20 - x^2)$ people per square mile, where
    $x$ is the distance from the seaside.\probonly{\footnote{The symbol
        $\rho$ is the Greek letter ``rho'' and is pronounced ``row'' (as
        in ``Row, row, row your boat'').}}  \textbf{Without using any
      calculus}, how could you \textbf{approximate} the population of the
    village?
  \end{problem}

  \begin{problemonly}
    If you come up with a method, here are some questions about it: 
    \begin{itemize}
    \item How accurate is your method?  If you needed a quick
      approximation that didn't require too much computation, could you
      come up with one?  (Try!)  What if you needed a really good
      approximation and were willing to spend a ton of time on
      computation? 
    \item How flexible is your method?  For example, if it was discovered
      in the next census that the density function $\rho(x)$ had changed
      to $500 (40 - \sqrt{x})$, would your method still work?
    \end{itemize}
  \end{problemonly}

  \begin{solution}
    The key is to try to use the simple situation of
    \ref{constant-density-warmup}: if a region has a \textit{constant}
    population density, we can find the population simply by multiplying
    this density by the area of the region.

    Unfortunately, the village in this problem does not have a constant
    population density.  But since we are only trying to approximate the
    population, here is one strategy we could use:
    \begin{enumerate}[label=\arabic*.]
    \item \textbf{Slice} the village into pieces so that, within each
      piece, the density is pretty close to being constant (that is, the
      density doesn't vary too much).
    \item Then, we can \textbf{approximate} the population of a single
      piece by multiplying the approximately constant density of that
      piece by the area of the piece.
    \item That gives us an approximate population for one piece of the
      village; if we can do that for all pieces of the village, we can
      simply \textbf{sum} these approximations to get the approximate
      population of the entire village.
    \end{enumerate}
    Let's follow this strategy.  First, we'll draw a picture of the
    problem so that we can figure out how to slice.  The village is a
    rectangle, and let's draw it so that the 6-inch sides are vertical and
    the sea lies on the left.  (These are totally arbitrary choices.)
    Since $x$ represents distance from the sea, it is equal to 0 on the
    side of the village that touches the sea and equal to 4 on the
    opposite side.  The left picture below shows the village, the sea, and
    the meaning of $x$:
    \begin{center}
      \includegraphics[height=1.5in]{images/density-rectangle-1}
      \hspace{0.5in}
      \includegraphics[height=1.5in]{images/density-rectangle-1-density}
    \end{center}
    The right picture shows a way to visualize the density: by color.  For
    instance, when $x = 0$, the density is $\rho(0) = 20000$ people per
    square mile (colored red on the picture); when $x = 4$, the density is
    $\rho(4) = 4000$ people per square mile (colored in blue).  In
    general, as we get further from the sea, the density decreases.

    Now, let's go through our strategy in more detail: 
    \begin{itemize}
    \item \textbf{Step 1: Slice.}

      We want to slice the village into pieces so that, within each piece,
      the density is pretty close to being constant.  As we can see from
      our pictures, the way to achieve this is by slicing vertically.  If
      we look at a thin vertical slice of the village, the density does
      not vary much.  (In contrast, if we look at a thin horizontal slice
      of the village, the density still varies a lot, between 20000 and
      4000 people per square mile.)

      For convenience, we'll slice the village into $n$ vertical strips of
      equal width.  Let's call this width $\Delta x$.  Since the entire
      width of the village is $4$ miles, $\Delta x = \frac{4}{n}$.  It's
      also useful to label the $x$-values where we make the slices, so
      we'll call these $x_0$, $x_1$, and so on, all the way up to $x_n$:
      \begin{center}
        \includegraphics[height=1.5in]{images/density-rectangle-1-slices}
      \end{center}
      Notice that $x_0 = 0$ and $x_n = 4$.  In general, $x_k$ is $k$
      ``slice widths'' to the right of $x_0$, so $x_k = x_0 + k \Delta x =
      k \Delta x$.

    \item \textbf{Step 2: Approximate the population within a single
        slice.}

      We need to approximate the population in each slice.  Rather than
      looking at the first slice, then the second, then the third, let's
      just look at a typical slice and think of it as being the ``$k$-th''
      slice.  Note that this slice is between $x = x_{k-1}$ and $x_k$ (for
      example, the first slice is between $x = x_0$ and $x = x_1$):
      \begin{center}
        \includegraphics[height=1.5in]{images/density-rectangle-1-kth-slice}
      \end{center}
      As we discussed in laying out our strategy, we will approximate the
      population in the $k$-th slice by saying
      \[ \textrm{population in $k$-th slice} \approx (\textrm{approximate
        population density in $k$-th slice}) \cdot (\textrm{area of $k$-th
        slice}).\] Since the population density does not vary much within
      a slice (by design --- that's how we sliced), we can just use the
      population density at any point in the slice as an approximation for
      the population density in the whole slice.  Let's use the population
      density at the right edge, which is $\rho(x_k)$.\footnote{There's no
        reason the right edge is better than the middle or the left edge
        or anywhere else, but it \textit{is} notationally convenient:
        we're talking about the $k$-th slice, and we can just use the
        density at a distance $x_k$ from the sea.}  The area of the $k$-th
      slice is pretty simple because the slice is just a rectangle: its
      height is 6 miles, and its width is $\Delta x$, so its area is $6
      \Delta x$.  So, we have:
      \begin{align*}
        \textrm{population density in $k$-th slice} & \approx \rho(x_k) \\
        \textrm{area of $k$-th slice} & = 6 \Delta x \\
        \shortintertext{Therefore,}
        \textrm{population in $k$-th slice} & \approx \rho(x_k) \cdot 6
        \Delta x
      \end{align*}

    \item \textbf{Step 3: Sum.}

      We figured out that the population in the $k$-th slice is
      approximately $\rho(x_k) \cdot 6 \Delta x$.  So, for instance, the
      population in the first slice ($k = 1$) is approximately $\rho(x_1)
      \cdot 6 \Delta x$, the population in the second slice ($k = 2$) is
      approximately $\rho(x_2) \cdot 6 \Delta x$, and so on.  To
      approximate the population of the whole village, we just add up the
      populations of all the slices.  This gives
      \[ \boxed{\rho(x_1) \cdot 6 \Delta x + \rho(x_2) \cdot 6 \Delta x +
        \cdots + \rho(x_n) \cdot 6 \Delta x}\] as our approximation for
      the population of the village.  We could also write this in
      summation notation as $\displaystyle \sum_{k=1}^n \rho(x_k) \cdot 6
      \Delta x$.  (Notice that this sum has $n$ terms, one per slice.)  If
      we wanted, we could also plug in the fact that $\rho(x) = 1000 (20 -
      x^2)$ to write the sum as
      \[ \boxed{\sum_{k=1}^n 1000 (20 - x_k^2) \cdot 6 \Delta x = 1000 (20
        - x_1^2) \cdot 6 \Delta x + \cdots + 1000 (20 - x_n^2) \cdot 6
        \Delta x}\]
    \end{itemize}
    There are a couple of things to notice about this strategy.  In
    general, if $n$ is really big, the approximation should be more
    accurate than if $n$ is small: in thinner slices, the density should
    vary less, which means the approximation should be more
    accurate.\footnote{This is a general trend, not an iron-clad rule.
      It's perfectly possible, depending on the density function, that
      using 10 slices will give you an approximation that happens to be
      closer to the actual population than the approximation you get using
      11 slices.  But, on the whole, as you increase the number of slices,
      the approximations should get better.}

    In addition, notice that we didn't really use the specific function
    $\rho(x) = 1000 (20 - x^2)$ much; if it had instead been $\rho(x) =
    500 (40 - \sqrt{x})$, we could have followed exactly the same
    procedure.  The important information was that density varies with
    distance to the sea, and that was really the only thing we used to
    decide which way to slice.
  \end{solution}
\end{enumerate}












\newpage







\end{document}
